\section{What establishes optimality}
\label{sec:verification}
A useful planner needs a feasible trajectory and a reason to trust its
reported performance. These are separate outputs. A small shooting
residual indicates a stationary candidate; it does not establish a
minimum. Comparing feasible branches improves an upper bound, while
a lower bound must cover every admissible competitor.

\subsection{The smooth second variation}
For an exact ideal extremal strictly outside the floor, let
$z=(\delta r,\delta v)$ and $\xi=\delta u$ satisfy
$z'=Az+B\xi$, $z(0)=z(h)=0$.
The fixed-endpoint second variation is
\begin{equation}
 Q[\xi]=\int_0^h\left(2|\xi|^2-z^TC_s z\right)dt,
 \label{eq:second}
\end{equation}
with $C_s=\operatorname{diag}(K,0)$ and
\begin{equation}
 \begin{split}
 K={}&\frac3{R^5}
 [p_vr^T+rp_v^T+(p_v\cdot r)I]\\
 &-\frac{15(p_v\cdot r)}{R^7}rr^T,\qquad R=|r|.
 \end{split}
\end{equation}
Differentiating $Dg=-I/R^3+3rr^T/R^5$ and contracting its
derivative with $p_v$ gives the displayed $K$.
The sign and factor two follow by augmenting the minimized cost with
$\int p\cdot(s'-f)dt$. Coercivity
$Q\ge\gamma_Q\norm{\xi}_2^2$ with $\gamma_Q>0$ proves a strict
local minimum in the $L^2$ control topology; a negative direction on
the endpoint kernel rules it out.

Initialize the vertical Jacobi frame $X(0)=0$, $Y(0)=I$ and integrate
\begin{equation}
 \begin{pmatrix}X\\Y\end{pmatrix}'
 =\begin{pmatrix}A&D\\-C_s&-A^T\end{pmatrix}
 \begin{pmatrix}X\\Y\end{pmatrix}.
\end{equation}
Conjugate times are rank losses of $X(t)$ at $t>0$.
For this regular fully actuated acceleration problem, absence of
conjugate points on $(0,h]$ gives strict local minimality; a first
conjugate point strictly before $h$ produces a negative direction
on the longer interval \cite{Bonnard2007}.
Where $X$ is invertible, $S=YX^{-1}$ satisfies
\begin{equation}
 S'+A^TS+SA+SDS+C_s=0,
\end{equation}
and completing the square yields
$Q=[z^TSz]_0^h+2\int|\xi-B^TSz/2|^2dt$.
Initialize the frame, not an arbitrary $S(0)=0$.

Sampling $\det X$ cannot exclude conjugate points: even multiplicity
can leave its sign unchanged. The supplied synthetic accessory problem
$Q/2=\int_0^1(z''^2-600z^2)dt$ in two identical coordinates
has such a crossing at $t_c=0.955712066275\ldots$.
It also has the explicit negative direction $z=1-\cos(2\pi t)$.
The test verifies the diagnostic logic, not a particular orbital branch.
Engine switches, radial contacts and flybys need their own junction
terms in a second variation.

\subsection{Global effort lower bounds}
With floor $a>0$, gravity obeys $|g|\le g_m=a^{-2}$.
For the endpoint moments $d_r,d_v$ of \eqref{eq:free},
\begin{equation}
 L_v=\frac{\pos{|d_v|-g_mh}^2}{h},\qquad
 L_r=\frac{3\pos{|d_r|-g_mh^2/2}^2}{h^3}
 \label{eq:momentbounds}
\end{equation}
are lower bounds. They follow by integrating \eqref{eq:kepler}
and applying Cauchy--Schwarz.
A joint directional bound, for $\lambda=(\alpha,\beta)\ne0$, is
\begin{equation}
 \begin{gathered}
 L_\lambda=
 \frac{\pos{\lambda\cdot(d_r,d_v)-g_m\int_0^h|q(t)|dt}^2}
 {\lambda^TW_0\lambda},\\
 q(t)=(h-t)\alpha+\beta.
 \end{gathered}
\end{equation}
An upper enclosure of the convex norm integral preserves this bound.
Several directions may be evaluated cheaply at the actual query.

Energy and angular momentum supply additional bounds for competitors
with $J\le U$, where $U\ge0$ is any cutoff. Put
\begin{equation}
 \begin{split}
 c_0&=|v_0|^2+2/a-2/|r_0|,\\
 I_v(U)&=[h\sqrt U+\sqrt{h^2U+hc_0}]^2,\\
 R_U&=|r_0|+h|v_0|+g_mh^2/2+\sqrt{Uh^3/3}.
 \end{split}
\end{equation}
To see the first estimate, write $I=\int_0^h|v|^2dt$.
The work identity and Cauchy--Schwarz imply
$|v(t)|^2\le c_0+2\sqrt{UI}$, whence $I\le hc_0+2h\sqrt{UI}$.
Solving this quadratic inequality in $\sqrt I$ gives $I_v(U)$.
Double integration of acceleration and Cauchy--Schwarz give $R_U$.
Equation~\eqref{eq:work} gives
$\int|v|^2dt\le I_v(U)$ and $|r|\le R_U$, hence
\begin{equation}
 L_E=\frac{(\Delta E)^2}{I_v(U)},\qquad
 L_M=\frac{(\Delta j)^2}{hR_U^2}.
\end{equation}
Degenerate zero numerators and denominators are handled by continuity.
A globally valid combination is
\begin{equation}
 L=\max\{0,L_v,L_r,L_\lambda,
                  \min[U,\max(L_E,L_M)]\}\le V.
 \label{eq:globalL}
\end{equation}
The cutoff clipping covers trajectories with effort greater than $U$.
An actually feasible trajectory of effort at most $U$ additionally
establishes $V\le U$ and suboptimality at most $U-L$.
Ideal lower bounds also apply to the bounded engine at the same
endpoints; its upper bound needs a band-feasible witness.

\subsection{A global certificate near a low-effort path}
Let $(\bar r,\bar v,\bar u)$ be an endpoint-exact feasible ideal
reference of effort $U$ and $p$ any continuous piecewise differentiable
multiplier. Define residuals
\begin{equation}
 e_u=2\bar u-p_v,\quad
 e_r=p_r'+Dg(\bar r)^Tp_v,\quad e_v=p_v'+p_r.
\end{equation}
For a competitor, write $\delta u=u-\bar u$,
$\delta r=r-\bar r$, $\delta v=v-\bar v$ and
$R_g=g(r)-g(\bar r)-Dg(\bar r)\delta r$.
Integration by parts gives the exact difference
\begin{equation}
 \begin{split}
 J-U={}&\norm{\delta u}_2^2+\int e_u\cdot\delta u\,dt\\
 &-\int(e_r\cdot\delta r+e_v\cdot\delta v+p_v\cdot R_g)dt.
 \end{split}
 \label{eq:gapidentity}
\end{equation}
Endpoint error would add a boundary term; it cannot be ignored.

Suppose $|\bar r|\ge b>a$, choose $0<d<b$ and put
$\ell=b-d$, $\omega=\sqrt{2/\ell^3}$.
Before a first exit from $|\delta r|<d$,
\begin{equation}
 |\delta r(t)|\le k_r(t)\norm{\delta u}_{2,[0,t]},\quad
 |\delta v(t)|\le k_v(t)\norm{\delta u}_{2,[0,t]},
\end{equation}
where
\begin{equation}
 k_r(t)^2=\int_0^t\frac{\sinh^2(\omega s)}{\omega^2}ds,\quad
 k_v(t)^2=\int_0^t\cosh^2(\omega s)ds.
\end{equation}
Every competitor with $J\le U$ has
$\norm{\delta u}_2\le2\sqrt U$. Thus
\begin{equation}
 2\sqrt U\,k_r(h)<d
 \label{eq:containment}
\end{equation}
prevents a first exit and covers all cheaper competitors.

Let $A_r=\int_0^h k_r^2dt$, $A_v=\int_0^h k_v^2dt$,
$P_v\ge\sup|p_v|$ and
\begin{equation}
 \begin{gathered}
 \gamma=1-\frac{3P_v}{\ell^4}A_r,\\
 \eta=\norm{e_u}_2+\sqrt{A_r}\norm{e_r}_2+
                    \sqrt{A_v}\norm{e_v}_2.
 \end{gathered}
\end{equation}
The sharp bound is $\norm{D^2g(x)[z,z]}\le6|z|^2/|x|^4$.
Indeed
\[
D^2g(x)[z,z]=\frac{6(x\cdot z)z+3|z|^2x}{|x|^5}
             -\frac{15(x\cdot z)^2x}{|x|^7}.
\]
For $|z|=1$ and $c=\hat x\cdot z$, its squared norm is
$9(1-2c^2+5c^4)/|x|^8\le36/|x|^8$, with equality radially.
Taylor's factor $1/2$ therefore gives $|R_g|\le3|\delta r|^2/\ell^4$.
This tightens the earlier factor 12 in $\gamma$ to 3 and implies $J-U\ge\gamma s^2-\eta s$ with $s=\norm{\delta u}_2$.
If \eqref{eq:containment} holds and $\gamma>0$,
\begin{equation}
 \boxed{\max(0,U-\eta^2/(4\gamma))\le V\le U.}
 \label{eq:tubegap}
\end{equation}
For an exact extremal $\eta=0$ proves uniqueness of the global
minimum, up to equality of controls almost everywhere. Numerical
use requires validated enclosures of the reference, residual norms
and constants. Sampled minima with padding do not supply that proof.

The earlier nonzero-effort fixture near $J=0.0032$ gives a small
double-precision estimate of the gap in \eqref{eq:tubegap}; it remains
an estimate, not a machine-verified certificate.
For difficult cells a verification function satisfying
$W_t+W_r\cdot v+W_v\cdot g-|W_v|^2/4\ge0$
throughout every possible cheaper trajectory yields another lower
bound by integration. A sampled residual mesh is insufficient.

\subsection{The atlas decision}
Store model, winding and active-set labels with every branch, together
with endpoint residuals, feasible reconstruction, effort enclosure,
and the domain and provenance of each optimality claim.
A point proof does not certify the surrounding interpolation cell.
Uniform continuation, constraint bounds and second-variation or
global-gap bounds must cover that cell.

At a query, combine valid lower bounds by their maximum and feasible
upper bounds by their minimum. Return infeasible if $L>\kappa$,
feasible if $U\le\kappa$, and unresolved otherwise.
A failed sufficient test is unresolved. An extra revolution alone is
not grounds for rejection: a circular coast over $4\pi$ is a two-turn
global zero-effort minimum for its own endpoints.
